\subsection{Caso 5. TCP: resets inesperados e indicios de inyección o desincronización}
\label{caso:tcp-resets}

\begin{situacion}
Las aplicaciones reportan sesiones TCP que terminan abruptamente o se reinician de forma repetida.
\end{situacion}

\paragraph{Procedimiento.}
\begin{itemize}[leftmargin=*, topsep=2pt, itemsep=2pt]
	\item En una PCAP de laboratorio, filtrar \texttt{tcp.flags.reset == 1} y correlacionar marcas de tiempo, endpoints y sesiones.
	\item Examinar si los RST coinciden con cierres legítimos, reinicios de servicio, firewall o pérdida de conectividad.
	\item Correlacionar con registros de Zeek/Suricata y registros de aplicación, si están instalados.
\end{itemize}

\paragraph{Indicadores que deben buscarse.}
RST fuera de contexto, frecuencia inusual, sesiones truncadas o inconsistencias que requieren correlación; no inferir MITM solo por observar RST.

\paragraph{Mitigación y defensa.}
Usar TLS correctamente validado; segmentar la red; actualizar los endpoints; activar detección IDS y revisar rutas y dispositivos intermedios.

\paragraph{Evidencias que se entregan.}
Tabla de flujos afectados, extractos de registros, tiempos y hash SHA-256 de la PCAP.

\paragraph{Preguntas de análisis.}
\begin{enumerate}[leftmargin=*, topsep=2pt, itemsep=2pt]
	\item ¿Qué causas benignas generan RST?
	\item ¿Qué evidencia adicional haría más sólida la hipótesis de manipulación?
\end{enumerate}

\registrohallazgo{Caso 5}
